\subsection{The six exceptional kernels and the remaining central sign}

\paragraph{D23.F26.}
Use $X,Y,Z=(XY)^{-1}$ and $S=SU(H)$ from D23.F25.  For the
six ordered pairs
\[
 (X,Y),\ (Y,Z),\ (Z,X),\ (Y,X),\ (Z,Y),\ (X,Z),
\tag{D23.F26.1}
\]
the simultaneous commutator-word evaluation is onto $S^6$.
Their six kernels on $F'$ are distinct even after quotienting
each target by its centre.  In particular the cyclic pair
$(Y,Z)$ does not have a swap-stable derived kernel.

This simultaneous map remains onto the direct product with
any prescribed finite word image whose order has prime divisors
only $2,3,7$.  Such finite images can encode all the old Krammer
pair values, their index-three induced evaluations of $C_f$,
and all needed finite residues of $b=\Psi_3/2$.

On the six exceptional values, impose the two-cycle, the
multiplier-one and multiplier-minus-one three-cycles, including
the correctly conjugated reversed cycles.  Whenever a tuple
satisfies these equations, independently multiplying two of
its three cyclic values by $-I$ gives another tuple satisfying
the same equations and all the same conjugation and generation
conditions.  In particular these finite conditions do not fix
the central sign of $f(X,Y)$ while the old data are held fixed.
They do not resolve the quadratic comparison solely through
the central factor $(-I)^b$ in the lifted orbital relation.
This conclusion does not discard the possible force of its
noncentral, projective periodic-product condition.

\begin{proof}
Every pair in (D23.F26.1) generates the same group
$H_0=\langle X,Y\rangle$, since any two of $X,Y,Z$ determine
the third.  Its derived subgroup is $S$ by D23.F25.  Write
$\overline S=S/\{I,-I\}$.  The conjugation-action orders of
the two marked generators on $S$ or $\overline S$ are,
in the six respective cases,
\[
       (36,6),\ (6,3),\ (3,36),\ (6,36),\ (3,6),\ (36,3).
\tag{D23.F26.2}
\]
These are all distinct.  Suppose two kernels of
$F'\to\overline S$ were equal.  The resulting identification
between their two quotient groups would intertwine conjugation
by each of the two free generators: both kernels are normal
in the full free group, and the quotient identification is
induced by the same words.  Thus the two ordered action orders
would agree.  Equation (D23.F26.2) rules this out.  This is a
kernel argument, not just a failure to find a linear intertwiner.
It also explains why the literal orders $6,6$ at $(Y,Z)$ are
misleading: $Z^3=-I$, whereas $Y^3$ is noncentral.

By D23.F25, $\overline S$ is simple.  A subdirect product of
copies of one nonabelian simple group is a product of diagonal
factors.  One can see the required statement inductively from
Goursat's lemma: a new simple factor is either independent or
is identified with a simple quotient of the preceding product;
an epimorphism from a product of simple groups to a simple
group factors through one coordinate, because the images of
the commuting factors are commuting normal subgroups.
Such an identification would equate two source kernels.
Their distinction therefore gives the full image
$\overline S^{\,6}$.

Let $J$ be the simultaneous image in $S^6$ and let
$Z_0=\{I,-I\}^6$.  The last paragraph gives $JZ_0=S^6$.
Because $Z_0$ is central, this makes $J$ normal and $S^6/J$
abelian.  But $S^6$ is perfect by D23.F25, so $J=S^6$.
This proves full independence including the six Schur-centre
coordinates, rather than only projective independence.

Now let $M$ be a finite word target with prime divisors among
$2,3,7$, and put $M_0=\operatorname{im}(F'\to M)$.
Every nontrivial quotient of $S$ is $S$ or $\overline S$:
a normal subgroup either lies in the centre, or surjects to
the simple projective quotient; in the latter case its quotient
in $S$ is abelian and hence trivial by perfectness.
Both nontrivial quotients have order divisible by $13$ and
$61$, as follows directly from D23.F25's order formula.
A nontrivial quotient of $S^6$ has a nontrivial image of some
factor $S$, so its order also has these prime divisors.
It cannot be a quotient of $M_0$.  Goursat's lemma proves
\[
 \operatorname{im}(F'\longrightarrow S^6\times M_0)
                         =S^6\times M_0.
\tag{D23.F26.3}
\]
Finite images of the ordinary free commutator subgroup and its
profinite closure agree, so all of these finite prescriptions
have common ordinary-word representatives.

Here is why the additional old data described in the statement
really fall under this assertion.  The old pair groups have
orders with prime divisors $2,3,7$ by D23.P6 and D23.C9;
the latter's perfect factors are $SU_3(3)$.  Given a finite
quotient $q:F\to M$ of this kind, induce it from the normal
index-three subgroup $F$ to $E=C_3\amalg C_3$.  The resulting
homomorphism takes values in $M\wr C_3$, which has no new
prime divisor.  Under the homomorphism from the parameter free
group to $E$ sending its two generators to $x,r$, the value
$C_f=f(x,r)$ lies in $F$, and its distinguished induced
coordinate is exactly $q(C_f)$.  Thus these $C_f$ values can
be encoded by ordinary finite homomorphisms of the parameter
free group, even if the original kernel of $q$ is not
$E$-stable.  Finally finite index-three Schreier coefficients
are encoded by the abelianization of the index-three subgroup
modulo the required modulus, followed by its finite coset
action.  To retain $b\bmod m$ use $\Psi_3\bmod2m$ and
restrict to its even values.  For the powers in the displayed
orbital matrices the required moduli have prime divisors
only $2,3$.  Finite products of all these targets remain of
the specified prime type.  No invariance of a congruence
kernel has been assumed in this induction argument.

We now verify the central-sign assertion at the level of all
the stated equations.  Denote the six values in order by
$v_0,\ldots,v_5$.  Two-cycle gives $v_{j+3}=v_j^{-1}$ for
$0\leq j\leq2$, and the two cycles at $(X,Y)$ are
\[
 v_2v_1v_0=I,\qquad
 v_2 Z^{-1}v_1Y^{-1}v_0X^{-1}=I.
\tag{D23.F26.4}
\]
For any $\alpha,\beta\in\mathbb F_2$ replace the first three
values by
\[
 v_0^*=(-I)^\alpha v_0,\qquad
 v_1^*=(-I)^\beta v_1,\qquad
 v_2^*=(-I)^{\alpha+\beta}v_2,
\tag{D23.F26.5}
\]
and replace their three reversed values by the inverses of
these new values.  Both products in (D23.F26.4) acquire the
same scalar $(-I)^{2\alpha+2\beta}=I$.
The reversed cycles use
$Yv_4Y^{-1},X^{-1}v_5X,v_3$, so the same verification applies
with the required inner corrections.  This proves preservation
of every cycle in the full outer three-puncture orbit.
Conjugation of a generator by a value $v_j$ is also unchanged
by its scalar multiplier.  Accordingly every finite test
depending only on the induced generator automorphism,
generation, or conjugacy classes of peripheral products is
unchanged.  Taking $\alpha=1$ changes the sign of $v_0$.
Equation (D23.F26.3) realizes the resulting tuple by one common
word while preserving all the encoded old data.

The exact literal orbital identity of D23.C16 contains
\[
 R\,A_E(x,V)\,Y^{-3b}\,f(X,Y)
       =A_E(x,y)\,Z_{123}^{\,b}(-I)^b,
\tag{D23.F26.6}
\]
where $Z_{123}=(K_1K_2)^3$; the letter $Z$ in
(D23.F26.1)--(D23.F26.4) instead denotes $(XY)^{-1}$.
We keep these two matrices explicitly distinct.
The last sign is essential: in the braid lift the two sides
without it have exponent sums $-6b$ and $6b$, so the exact
central correction is $z_4^{-b}$; the Krammer image of $z_4$
is $-I$.  At fixed old data, the scalar change (D23.F26.5)
can select either lift of the required projective value of
$f(X,Y)$ without affecting the other finite conditions just
proved.  Thus once such a projective value has a solution of
the literal cycle conditions, its sign in (D23.F26.6) supplies
no further obstruction on the old data.  This precise assertion
does not assert that its noncentral value always satisfies
the periodic-product condition, nor that one common finite
word is an actual gentle or full-GT parameter.

The exact action orders and scalar powers used here are part
of \path{agents/five_level_six/krammer_orbital_unitary.py} and
its JSON record.  The central-sign verification is an identity
in any group with central involution and is written out in
(D23.F26.4)--(D23.F26.5); it requires no numerical approximation.
\end{proof}
